% Kooktemperaturen: a \subchapter-level (H2) heading grouped under the
% single "Bijlagen" chapter-level heading in main.tex.
\subchapter{Kooktemperaturen}

\lettrine{B}{ij} een aantal vaste temperaturen verandert eten van structuur,
van eiwit dat stolt tot suiker die bruin wordt. Hieronder staat, van koud
naar heet, wat er ongeveer gebeurt.

% Verticale tijdlijn (4-230+ °C): zelfde opzet als de matenlijn in Kleine
% maten, maar hier telkens één temperatuur in plaats van een gewichtsbereik.
\begin{center}
\begin{tikzpicture}[y=1cm]
\draw[color=terracotta, line width=0.6pt, {Stealth[length=1.8mm,width=1.1mm]}-{Stealth[length=1.8mm,width=1.1mm]}] (0,0.3) -- (0,-13.05);

\node at (0,0) {\textcolor{terracotta}{\scriptsize$\blacklozenge$}};
\draw[color=inkfaint, line width=0.3pt] (0,0) -- (-0.5,0);
\node[align=right, text width=4.8cm, anchor=east] at (-0.6,0)
  {\textbf{\textcolor{terracotta}{4\,°C}}\\[0pt]
   {\footnotesize\textcolor{inkfaint}{Enzymen vertragen, bederf gaat trager.}}};

\node at (0,-0.75) {\textcolor{terracotta}{\scriptsize$\blacklozenge$}};
\draw[color=inkfaint, line width=0.3pt] (0,-0.75) -- (0.5,-0.75);
\node[align=left, text width=4.8cm, anchor=west] at (0.6,-0.75)
  {\textbf{\textcolor{terracotta}{15\,°C}}\\[0pt]
   {\footnotesize\textcolor{inkfaint}{Boter wordt zacht, aroma's komen los.}}};

\node at (0,-1.50) {\textcolor{terracotta}{\scriptsize$\blacklozenge$}};
\draw[color=inkfaint, line width=0.3pt] (0,-1.50) -- (-0.5,-1.50);
\node[align=right, text width=4.8cm, anchor=east] at (-0.6,-1.50)
  {\textbf{\textcolor{terracotta}{34\,°C}}\\[0pt]
   {\footnotesize\textcolor{inkfaint}{Cacaoboter in chocolade smelt.}}};

\node at (0,-2.25) {\textcolor{terracotta}{\scriptsize$\blacklozenge$}};
\draw[color=inkfaint, line width=0.3pt] (0,-2.25) -- (0.5,-2.25);
\node[align=left, text width=4.8cm, anchor=west] at (0.6,-2.25)
  {\textbf{\textcolor{terracotta}{40\,°C}}\\[0pt]
   {\footnotesize\textcolor{inkfaint}{Vis-eiwit stolt, verliest zijn glazigheid.}}};

\node at (0,-3.00) {\textcolor{terracotta}{\scriptsize$\blacklozenge$}};
\draw[color=inkfaint, line width=0.3pt] (0,-3.00) -- (-0.5,-3.00);
\node[align=right, text width=4.8cm, anchor=east] at (-0.6,-3.00)
  {\textbf{\textcolor{terracotta}{50\,°C}}\\[0pt]
   {\footnotesize\textcolor{inkfaint}{Vlees-eiwit begint te stollen.}}};

\node at (0,-3.75) {\textcolor{terracotta}{\scriptsize$\blacklozenge$}};
\draw[color=inkfaint, line width=0.3pt] (0,-3.75) -- (0.5,-3.75);
\node[align=left, text width=4.8cm, anchor=west] at (0.6,-3.75)
  {\textbf{\textcolor{terracotta}{60\,°C}}\\[0pt]
   {\footnotesize\textcolor{inkfaint}{Eiwit van een ei stolt en wordt wit.}}};

\node at (0,-4.50) {\textcolor{terracotta}{\scriptsize$\blacklozenge$}};
\draw[color=inkfaint, line width=0.3pt] (0,-4.50) -- (-0.5,-4.50);
\node[align=right, text width=4.8cm, anchor=east] at (-0.6,-4.50)
  {\textbf{\textcolor{terracotta}{65\,°C}}\\[0pt]
   {\footnotesize\textcolor{inkfaint}{Eigeel stolt, zetmeel bindt vocht.}}};

\node at (0,-5.25) {\textcolor{terracotta}{\scriptsize$\blacklozenge$}};
\draw[color=inkfaint, line width=0.3pt] (0,-5.25) -- (0.5,-5.25);
\node[align=left, text width=4.8cm, anchor=west] at (0.6,-5.25)
  {\textbf{\textcolor{terracotta}{70\,°C}}\\[0pt]
   {\footnotesize\textcolor{inkfaint}{Taai bindweefsel wordt zachte gelatine.}}};

\node at (0,-6.00) {\textcolor{terracotta}{\scriptsize$\blacklozenge$}};
\draw[color=inkfaint, line width=0.3pt] (0,-6.00) -- (-0.5,-6.00);
\node[align=right, text width=4.8cm, anchor=east] at (-0.6,-6.00)
  {\textbf{\textcolor{terracotta}{80\,°C}}\\[0pt]
   {\footnotesize\textcolor{inkfaint}{Vrijwel alle eiwit is nu gestold.}}};

\node at (0,-6.75) {\textcolor{terracotta}{\scriptsize$\blacklozenge$}};
\draw[color=inkfaint, line width=0.3pt] (0,-6.75) -- (0.5,-6.75);
\node[align=left, text width=4.8cm, anchor=west] at (0.6,-6.75)
  {\textbf{\textcolor{terracotta}{85\,°C}}\\[0pt]
   {\footnotesize\textcolor{inkfaint}{Pectine breekt af, groente wordt snel zacht.}}};

\node at (0,-7.50) {\textcolor{terracotta}{\scriptsize$\blacklozenge$}};
\draw[color=inkfaint, line width=0.3pt] (0,-7.50) -- (-0.5,-7.50);
\node[align=right, text width=4.8cm, anchor=east] at (-0.6,-7.50)
  {\textbf{\textcolor{terracotta}{95\,°C}}\\[0pt]
   {\footnotesize\textcolor{inkfaint}{Water suddert, net onder de kook.}}};

\node at (0,-8.25) {\textcolor{terracotta}{\scriptsize$\blacklozenge$}};
\draw[color=inkfaint, line width=0.3pt] (0,-8.25) -- (0.5,-8.25);
\node[align=left, text width=4.8cm, anchor=west] at (0.6,-8.25)
  {\textbf{\textcolor{terracotta}{100\,°C}}\\[0pt]
   {\footnotesize\textcolor{inkfaint}{Water kookt.}}};

\node at (0,-9.00) {\textcolor{terracotta}{\scriptsize$\blacklozenge$}};
\draw[color=inkfaint, line width=0.3pt] (0,-9.00) -- (-0.5,-9.00);
\node[align=right, text width=4.8cm, anchor=east] at (-0.6,-9.00)
  {\textbf{\textcolor{terracotta}{120\,°C}}\\[0pt]
   {\footnotesize\textcolor{inkfaint}{Maillard-reactie begint, zodra het oppervlak droog is: het begin van bruinen.}}};

\node at (0,-9.75) {\textcolor{terracotta}{\scriptsize$\blacklozenge$}};
\draw[color=inkfaint, line width=0.3pt] (0,-9.75) -- (0.5,-9.75);
\node[align=left, text width=4.8cm, anchor=west] at (0.6,-9.75)
  {\textbf{\textcolor{terracotta}{140\,°C}}\\[0pt]
   {\footnotesize\textcolor{inkfaint}{Maillard-reactie op zijn sterkst.}}};

\node at (0,-10.50) {\textcolor{terracotta}{\scriptsize$\blacklozenge$}};
\draw[color=inkfaint, line width=0.3pt] (0,-10.50) -- (-0.5,-10.50);
\node[align=right, text width=4.8cm, anchor=east] at (-0.6,-10.50)
  {\textbf{\textcolor{terracotta}{160\,°C}}\\[0pt]
   {\footnotesize\textcolor{inkfaint}{Karamellisatie van suiker begint.}}};

\node at (0,-11.25) {\textcolor{terracotta}{\scriptsize$\blacklozenge$}};
\draw[color=inkfaint, line width=0.3pt] (0,-11.25) -- (0.5,-11.25);
\node[align=left, text width=4.8cm, anchor=west] at (0.6,-11.25)
  {\textbf{\textcolor{terracotta}{180\,°C}}\\[0pt]
   {\footnotesize\textcolor{inkfaint}{Risico op acrylamide bij hard gebakken zetmeel.}}};

\node at (0,-12.00) {\textcolor{terracotta}{\scriptsize$\blacklozenge$}};
\draw[color=inkfaint, line width=0.3pt] (0,-12.00) -- (-0.5,-12.00);
\node[align=right, text width=4.8cm, anchor=east] at (-0.6,-12.00)
  {\textbf{\textcolor{terracotta}{200\,°C}}\\[0pt]
   {\footnotesize\textcolor{inkfaint}{Rookpunt van de meeste bakolie.}}};

\node at (0,-12.75) {\textcolor{terracotta}{\scriptsize$\blacklozenge$}};
\draw[color=inkfaint, line width=0.3pt] (0,-12.75) -- (0.5,-12.75);
\node[align=left, text width=4.8cm, anchor=west] at (0.6,-12.75)
  {\textbf{\textcolor{terracotta}{Vanaf 230\,°C}}\\[0pt]
   {\footnotesize\textcolor{inkfaint}{Eten verkoolt en wordt bitter.}}};
\end{tikzpicture}
\end{center}

\vspace{0.6em}
{\footnotesize\itshape\color{inkfaint}Richtwaarden uit de
keukenscheikunde\cite{mcgee-onfoodandcooking}, geen exacte omslagpunten: het
echte moment hangt af van het ingrediënt zelf.\par}
